\documentclass[11pt,a4paper]{article}
\usepackage[T1]{fontenc}
\usepackage{lmodern}
\usepackage{amsmath,amssymb,mathtools,booktabs,array}
\usepackage[margin=24mm,headheight=14pt]{geometry}
\usepackage{microtype,xcolor,hyperref,fancyhdr,tcolorbox,enumitem}
\hypersetup{colorlinks=true,linkcolor=blue!45!black,urlcolor=blue!45!black,pdftitle={From Geometry to the Einstein--Cartan Action}}
\definecolor{deep}{HTML}{17344D}\definecolor{soft}{HTML}{F1F5F8}
\pagestyle{fancy}\fancyhf{}\fancyhead[L]{\small \textsc{Mathematical foundations of Einstein--Cartan gravity}}\fancyhead[R]{\small Chapter 1}\fancyfoot[C]{\thepage}
\setlength{\parindent}{0pt}\setlength{\parskip}{6pt}
\setcounter{secnumdepth}{2}
\newtcolorbox{whybox}{colback=soft,colframe=deep!40,boxrule=.6pt,arc=2mm}
\newcommand{\LC}{\mathring{\nabla}}
\newcommand{\C}{\left\{\begin{smallmatrix}\rho\\\mu\nu\end{smallmatrix}\right\}}
\newcommand{\kappaG}{\kappa}
\newcommand{\RR}{\mathring{R}}
\begin{document}
\begin{center}
{\small EINSTEIN--CARTAN RESEARCH PROJECT \quad / \quad MATHEMATICAL DERIVATIONS}\par\vspace{5mm}
{\LARGE\bfseries\color{deep} Chapter 1\\ From Geometry to the Einstein--Cartan Action}\par\vspace{2mm}
{\large A first-principles textbook chapter, before field variation}\par\vspace{3mm}
{\small 8 November 2026 \quad $\bullet$ \quad Notation: $(-,+,+,+)$, $c=1$}\par
\end{center}
\vspace{2mm}
\begin{whybox}
\textbf{Reading goal.} By the end of this chapter you should be able to explain \emph{why} we choose an action for gravity, \emph{how} the metric and connection define torsion and curvature, and \emph{how} replacing the Levi--Civita connection by a metric-compatible Einstein--Cartan connection produces additional contortion terms in the gravitational action. We deliberately \textbf{stop before varying the action}. No Einstein or Cartan field equation is derived here.
\end{whybox}
\tableofcontents
\newpage
\section{From classical action to gravitational action}
\subsection{What is an action, and why use one?}
In analytical mechanics one introduces a path $q(t)$ and an action
\begin{equation}\label{eq:particle}
S[q]=\int_{t_1}^{t_2}L(q,\dot q,t)\,dt.
\end{equation}
For $L=\tfrac12m\dot q^2-V(q)$, requiring $\delta S=0$ for variations with fixed endpoints yields the Euler--Lagrange equation $d(\partial L/\partial\dot q)/dt-\partial L/\partial q=0$, and hence $m\ddot q=-\partial_qV$.

The action is therefore an organized way of specifying \emph{dynamics}, not merely describing a configuration. A metric tells us the geometry of spacetime; an action tells us which possible metrics satisfy the theory's equations.

A field $\phi(x)$ has values at every spacetime point. Its action is an integral over four coordinates, schematically $S[\phi]=\int d^4x\,\mathcal L(\phi,\partial\phi)$. We will apply this idea to gravitational variables.

\subsection{Why the Ricci scalar?}
In metric general relativity the fundamental geometric field is the symmetric metric $g_{\mu\nu}(x)$. We seek a local, coordinate-invariant gravitational action. A natural starting assumption is to build it from the metric and a small number of its derivatives, and to choose a minimal form giving second-order metric equations after treating boundary terms correctly.

The curvature of the metric's torsion-free connection produces a scalar $R(g)$, called the \emph{Ricci scalar}. This is an appropriate candidate because the integrand of a coordinate-invariant action must transform as a scalar density. However, $R$ is not a uniquely forced choice: cosmological-constant terms, higher-curvature terms, and alternative gravitational theories are possible. The Einstein--Hilbert choice is the simplest standard curvature term with the required low-energy properties.

\subsection{Why is there a factor $\sqrt{-g}$?}
A coordinate volume $d^4x=dt\,dx\,dy\,dz$ changes when we replace one set of coordinates by another. The determinant $g=\det(g_{\mu\nu})$ transforms so that
\begin{equation}\label{eq:volume}
dV_4=\sqrt{-g}\,d^4x
\end{equation}
is invariant under orientation-preserving coordinate changes (for a Lorentzian metric of signature $-+++$). Thus, $\sqrt{-g}R\,d^4x$ can be integrated over spacetime without referring to any privileged coordinates.

\subsection{Why $1/(2\kappa)$?}
Before determining its normalization, we can write
\begin{equation}
S_g=A\int d^4x\,\sqrt{-g}\,R(g).
\end{equation}
The coefficient $A$ is fixed by matching the resulting gravitational field equation to Newtonian gravity in the weak-field, slow-motion limit. With $c=1$, this gives $A=1/(16\pi G)=1/(2\kappa)$, where $\kappa=8\pi G$. For orientation, the weak-field relation $g_{00}\simeq-(1+2\Phi)$, together with the linearized $00$ Einstein equation, gives $2\nabla^2\Phi=\kappa\rho$, reproducing $\nabla^2\Phi=4\pi G\rho$. We are \emph{motivating} the coefficient here, not deriving the Einstein field equations in this chapter.

The standard metric gravitational action is therefore
\begin{equation}\label{eq:EH}
\boxed{S_{\mathrm{EH}}[g]=\frac1{2\kappa}\int_M d^4x\,\sqrt{-g}\,R(g).}
\end{equation}
A matter action $S_m[g,\psi]$ is added to describe the non-gravitational fields $\psi$. A well-posed metric variation on a manifold with a boundary also requires appropriate boundary terms, such as the Gibbons--Hawking--York term for standard Dirichlet boundary conditions. We keep boundary details explicit when needed but focus on bulk expressions.

\begin{whybox}\textbf{Why this matters next.} Equation \eqref{eq:EH} is the baseline that we extend. In Einstein--Cartan theory, we retain the curvature-based action principle but replace the restricted torsion-free connection by a more general \emph{metric-compatible} connection. Its curvature scalar is then $\widetilde R(g,K)$ instead of $R(g)$.\end{whybox}

\section{Metric, connection and parallel transport}
\subsection{The metric measures intervals}
For nearby events $x^\mu$ and $x^\mu+dx^\mu$, the metric defines the infinitesimal interval
\begin{equation}\label{eq:metric}
\boxed{ds^2=g_{\mu\nu}(x)\,dx^\mu dx^\nu.}
\end{equation}
The inverse metric $g^{\mu\nu}$ satisfies $g^{\mu\alpha}g_{\alpha\nu}=\delta^\mu_\nu$, where $\delta^\mu_\nu$ is the Kronecker delta. We raise and lower indices with the metric. Repeated upper/lower indices are summed (Einstein convention).

For example, a flat plane written in polar coordinates satisfies $d\ell^2=dr^2+r^2d\varphi^2$. The coefficient $r^2$ reflects the choice of coordinates; the plane remains flat. \emph{A nonconstant metric coefficient does not automatically mean nonzero curvature.}

\subsection{Why do we need a connection?}
The metric measures intervals, but does not alone specify a general arbitrary rule for comparing vector components at different points. That comparison is expressed through a connection. In this chapter we fix the following index convention:
\begin{equation}\label{eq:convention}
\boxed{\nabla_\mu V^\rho=\partial_\mu V^\rho+\Gamma^\rho{}_{\mu\nu}V^\nu.}
\end{equation}
Thus, the \emph{first lower} connection index labels the differentiation direction. This ordering must remain fixed throughout every derivation. Some textbooks use the opposite ordering for the lower indices.

Along a curve $x^\mu(\lambda)$, a parallel-transported vector satisfies
\begin{equation}
\frac{dV^\rho}{d\lambda}+\Gamma^\rho{}_{\mu\nu}\frac{dx^\mu}{d\lambda}V^\nu=0.
\end{equation}
Even in flat polar coordinates, the \emph{components} of a parallel-transported vector can change because the coordinate basis changes. Curvature cannot be inferred merely from such component changes.

\subsection{Levi--Civita connection and metric compatibility}
In ordinary metric GR the connection is chosen to be (i) metric-compatible, $\nabla_\lambda g_{\mu\nu}=0$, and (ii) torsion-free. These conditions determine the unique Levi--Civita connection:
\begin{equation}\label{eq:LC}
\boxed{\left\{^\rho{}_{\mu\nu}\right\}=\tfrac12g^{\rho\sigma}(\partial_\mu g_{\sigma\nu}+\partial_\nu g_{\sigma\mu}-\partial_\sigma g_{\mu\nu}).}
\end{equation}
Metric compatibility means that parallel transport preserves metric inner products. It does \emph{not} require torsion to vanish. Einstein--Cartan theory keeps metric compatibility but removes the torsion-free restriction.

\section{Torsion and contortion from the connection}
\subsection{Definition of torsion}
With our convention, the torsion tensor is the antisymmetric part of the connection in its two lower indices:
\begin{equation}\label{eq:torsion}
\boxed{T^\rho{}_{\mu\nu}=\widetilde\Gamma^\rho{}_{\mu\nu}-\widetilde\Gamma^\rho{}_{\nu\mu}=2\widetilde\Gamma^\rho{}_{[\mu\nu]}.}
\end{equation}
The notation $A_{[\mu\nu]}=\tfrac12(A_{\mu\nu}-A_{\nu\mu})$ means antisymmetrization. Unlike the connection coefficients themselves, torsion is a tensor. The factor of $2$ is conventional: for example, Br\'echet et al. use $Q^\rho{}_{\mu\nu}=\widetilde\Gamma^\rho{}_{[\mu\nu]}=T^\rho{}_{\mu\nu}/2$.

\subsection{Why split the connection?}
Under metric compatibility we may compare the full connection with the unique Levi--Civita connection of the same metric:
\begin{equation}\label{eq:split}
\boxed{\widetilde\Gamma^\rho{}_{\mu\nu}=\left\{^\rho{}_{\mu\nu}\right\}+K^\rho{}_{\mu\nu}.}
\end{equation}
The difference $K^\rho{}_{\mu\nu}=\widetilde\Gamma^\rho{}_{\mu\nu}-\{^\rho{}_{\mu\nu}\}$ is the \emph{contortion tensor}. This is an algebraic decomposition, not an additional physical field equation. There is no circular definition: once $g$ and torsion are specified, metric compatibility determines $K$ and hence the full connection.

\subsection{Derive contortion in terms of torsion}
Lower the first index by $K_{\rho\mu\nu}=g_{\rho\lambda}K^\lambda{}_{\mu\nu}$. Because both full and Levi--Civita connections are metric-compatible, their difference obeys
\begin{equation}\label{eq:antiK}
K_{\rho\mu\nu}=-K_{\nu\mu\rho}.
\end{equation}
The torsion-free Levi--Civita terms cancel under antisymmetrization, so
\begin{equation}
T_{\rho\mu\nu}=K_{\rho\mu\nu}-K_{\rho\nu\mu}.
\end{equation}
Using \eqref{eq:antiK}, rewrite the second term: $K_{\rho\nu\mu}=-K_{\mu\nu\rho}$. Thus
\begin{align*}
T_{\rho\mu\nu}&=K_{\rho\mu\nu}+K_{\mu\nu\rho},\\
T_{\mu\nu\rho}&=K_{\mu\nu\rho}+K_{\nu\rho\mu},\\
T_{\nu\rho\mu}&=K_{\nu\rho\mu}+K_{\rho\mu\nu}.
\end{align*}
Add the first and third relations, and subtract the second. The undesired components cancel, leaving
\begin{equation}\label{eq:contortion}
\boxed{K_{\rho\mu\nu}=\tfrac12\left(T_{\rho\mu\nu}-T_{\mu\nu\rho}+T_{\nu\rho\mu}\right).}
\end{equation}
This identity is specific to our stated index convention. Contortion signs in the literature may look opposite because authors sometimes define $\widetilde\Gamma=\{\}-K$ rather than \eqref{eq:split}.

\begin{whybox}\textbf{Physical versus mathematical statement.} Equations \eqref{eq:torsion}--\eqref{eq:contortion} define geometric quantities and relate them algebraically. They do \emph{not} yet say that spin causes torsion. That statement follows later from the Cartan field equation, derived by varying an action.\end{whybox}

\section{How the connection produces curvature}
\subsection{From the commutator to the Riemann tensor}
A connection allows two successive covariant derivatives. Their commutator measures a noncommutativity associated with curvature, plus an additional term when torsion is present. With our convention,
\begin{equation}\label{eq:commutator}
[\nabla_\mu,\nabla_\nu]V^\rho=\widetilde R^\rho{}_{\sigma\mu\nu}V^\sigma-T^\lambda{}_{\mu\nu}\nabla_\lambda V^\rho.
\end{equation}
The curvature tensor appearing here is constructed by differentiating and multiplying connection coefficients. In a coordinate basis it is
\begin{equation}\label{eq:riemann}
\boxed{\begin{aligned}
\widetilde R^\rho{}_{\sigma\mu\nu}={}&\partial_\mu\widetilde\Gamma^\rho{}_{\nu\sigma}-\partial_\nu\widetilde\Gamma^\rho{}_{\mu\sigma}\\
&+\widetilde\Gamma^\rho{}_{\mu\lambda}\widetilde\Gamma^\lambda{}_{\nu\sigma}-\widetilde\Gamma^\rho{}_{\nu\lambda}\widetilde\Gamma^\lambda{}_{\mu\sigma}.
\end{aligned}}
\end{equation}
To see \emph{why} there are products of two connections: applying $\nabla_\mu$ to $\nabla_\nu V$ differentiates both $V$ and the coefficients $\Gamma$; the derivative of $V$ already contains a connection term. The combined products of two $\Gamma$s are the resulting quadratic contributions. Equation \eqref{eq:riemann} specifies our Riemann-index and sign convention; alternate textbooks may choose the opposite sign.

\subsection{Ricci tensor and Ricci scalar}
The Riemann tensor has four indices. Contracting the first and third gives the \emph{Ricci tensor}:
\begin{equation}\label{eq:ricci}
\boxed{\widetilde R_{\sigma\nu}=\widetilde R^\rho{}_{\sigma\rho\nu}.}
\end{equation}
A further contraction with the inverse metric gives the \emph{Ricci scalar}:
\begin{equation}\label{eq:ricciscalar}
\boxed{\widetilde R=g^{\sigma\nu}\widetilde R_{\sigma\nu}.}
\end{equation}
In a torsionful geometry, the Ricci tensor need not be symmetric, although its antisymmetric part vanishes when contracted with the symmetric metric. The scalar in \eqref{eq:ricciscalar} is the curvature quantity entering the minimal Einstein--Cartan gravitational action.

\section{The exact curvature decomposition using contortion}
\subsection{Insert $\widetilde\Gamma=\{\}+K$}
We now use the precise reason for introducing contortion: it lets us express Einstein--Cartan curvature in terms of familiar GR curvature plus extra terms. In each occurrence of $\widetilde\Gamma$ in \eqref{eq:riemann}, substitute equation \eqref{eq:split}.

The terms containing only Levi--Civita connections form $R^\rho{}_{\sigma\mu\nu}(\{\})$. The pieces linear in $K$ contain $\partial K$ and cross-products $\{\}K$; these recombine into the \emph{Levi--Civita covariant derivative} $\LC_\mu K$. The remaining pieces contain $KK$. This gives the exact tensor identity
\begin{equation}\label{eq:Rsplit}
\boxed{\begin{aligned}
\widetilde R^\rho{}_{\sigma\mu\nu}={}&\RR^\rho{}_{\sigma\mu\nu}
+\LC_\mu K^\rho{}_{\nu\sigma}-\LC_\nu K^\rho{}_{\mu\sigma}\\
&+K^\rho{}_{\mu\lambda}K^\lambda{}_{\nu\sigma}-K^\rho{}_{\nu\lambda}K^\lambda{}_{\mu\sigma}.
\end{aligned}}
\end{equation}
Here $\LC$ and $\RR$ always mean the torsion-free Levi--Civita derivative and curvature of the same metric $g$. They are not new independent fields.

\subsection{Check the covariant derivative terms}
To make the middle step transparent, for a $(1,2)$ tensor the Levi--Civita derivative is
\begin{align}\label{eq:nablaK}
\LC_\mu K^\rho{}_{\nu\sigma}
={}&\partial_\mu K^\rho{}_{\nu\sigma}+\left\{^\rho{}_{\mu\lambda}\right\}K^\lambda{}_{\nu\sigma}\notag\\
&-\left\{^\lambda{}_{\mu\nu}\right\}K^\rho{}_{\lambda\sigma}-\left\{^\lambda{}_{\mu\sigma}\right\}K^\rho{}_{\nu\lambda}.
\end{align}
Subtract the expression with $\mu\leftrightarrow\nu$. Because $\{^\lambda{}_{\mu\nu}\}=\{^\lambda{}_{\nu\mu}\}$, the terms multiplying $K^\rho{}_{\lambda\sigma}$ cancel. The other terms reproduce the cross-products $\{\}K$ from the direct expansion. This is why the shorthand $R(\widetilde\Gamma)=R(\{\})+\LC K+K^2$ is legitimate, even though it suppresses important index contractions.

\subsection{Contract to obtain the Ricci tensor}
Set $\mu=\rho$ and sum in \eqref{eq:Rsplit}. Then
\begin{equation}\label{eq:Riccisplit}
\boxed{\begin{aligned}
\widetilde R_{\sigma\nu}={}&\RR_{\sigma\nu}
+\LC_\rho K^\rho{}_{\nu\sigma}-\LC_\nu K^\rho{}_{\rho\sigma}\\
&+K^\rho{}_{\rho\lambda}K^\lambda{}_{\nu\sigma}
-K^\rho{}_{\nu\lambda}K^\lambda{}_{\rho\sigma}.
\end{aligned}}
\end{equation}
Notice that the contraction does not simply \emph{remove} an index: it sets two indices equal and sums over their four possible values.

\subsection{Contract again: the complete Ricci scalar}
Multiply by $g^{\sigma\nu}$. This produces the full scalar-curvature decomposition:
\begin{equation}\label{eq:scalardecomp}
\boxed{\begin{aligned}
\widetilde R={}&\RR
+g^{\sigma\nu}\left(\LC_\rho K^\rho{}_{\nu\sigma}-\LC_\nu K^\rho{}_{\rho\sigma}\right)\\
&+g^{\sigma\nu}\left(K^\rho{}_{\rho\lambda}K^\lambda{}_{\nu\sigma}
-K^\rho{}_{\nu\lambda}K^\lambda{}_{\rho\sigma}\right).
\end{aligned}}
\end{equation}
\emph{This equation is the main mathematical result of Chapter 1.} It is exact for the specified connection decomposition and curvature convention. No field equations have been used.

\subsection{Why the derivative contribution becomes a boundary term}
Since the Levi--Civita derivative preserves the metric, $\LC_\rho g^{\mu\nu}=0$, we can move the metric through the derivative in equation \eqref{eq:scalardecomp}. Define two vectors
\begin{equation}
A^\rho=g^{\sigma\nu}K^\rho{}_{\nu\sigma},\qquad
B^\nu=g^{\sigma\nu}K^\rho{}_{\rho\sigma}.
\end{equation}
Then the derivative terms combine exactly into
\begin{equation}\label{eq:div}
g^{\sigma\nu}\left(\LC_\rho K^\rho{}_{\nu\sigma}-\LC_\nu K^\rho{}_{\rho\sigma}\right)
=\LC_\mu(A^\mu-B^\mu).
\end{equation}
For a Levi--Civita divergence,
\begin{equation}\label{eq:dividentity}
\boxed{\sqrt{-g}\,\LC_\mu V^\mu=\partial_\mu(\sqrt{-g}\,V^\mu).}
\end{equation}
Therefore its spacetime integral becomes a boundary integral by the divergence theorem. Under appropriate boundary conditions, removing it does not alter the \emph{bulk} Euler--Lagrange equations. This statement does not mean the boundary is physically irrelevant; a complete variational principle must keep track of suitable boundary contributions.

The remaining quadratic expression has a precise definition:
\begin{equation}\label{eq:QK}
\boxed{Q(K)=g^{\sigma\nu}\left(K^\rho{}_{\rho\lambda}K^\lambda{}_{\nu\sigma}
-K^\rho{}_{\nu\lambda}K^\lambda{}_{\rho\sigma}\right).}
\end{equation}
It is \emph{not} a generic $K^2$ guessed by dimensional analysis: it is a specific contraction derived from the Riemann curvature formula.

\begin{whybox}\textbf{Consistency check.} Setting $K^\rho{}_{\mu\nu}=0$ removes both the derivative and quadratic terms; then $\widetilde R=\RR$, precisely the torsion-free Ricci scalar. This is the correct GR limit.\end{whybox}

\section{The minimal Einstein--Cartan action: our stopping point}
\subsection{What changes relative to Einstein--Hilbert gravity?}
We retain the simplest curvature-based gravitational action but replace the Levi--Civita scalar with the scalar of the full metric-compatible connection. Denote collectively by $\psi$ the matter fields, which may carry intrinsic spin. The action is
\begin{equation}\label{eq:EC}
\boxed{S[g,K,\psi]=\frac{1}{2\kappa}\int_M d^4x\,\sqrt{-g}\,\widetilde R(g,K)+S_m[g,K,\psi].}
\end{equation}
The full connection is given by $\widetilde\Gamma=\{g\}+K$, so the gravitational variables may be taken to be the metric and contortion. Independent variations of the metric and of the \emph{metric-compatible} contortion are equivalent to choosing the metric and a corresponding metric-compatible connection as variables, after the appropriate field redefinition and matter treatment. This choice avoids introducing tetrads at the outset, although tetrads and a spin connection are particularly natural for fundamental spinor fields.

\subsection{Insert the exact decomposition}
Using equations \eqref{eq:scalardecomp}--\eqref{eq:QK}, and displaying separately the boundary contribution $S_{\partial M}$, we obtain
\begin{equation}\label{eq:ECbulk}
\boxed{\begin{aligned}
S[g,K,\psi]={}&\frac{1}{2\kappa}\int_M d^4x\,\sqrt{-g}\,\left[\RR(g)+Q(K)\right]\\
&+S_m[g,K,\psi]+S_{\partial M}.
\end{aligned}}
\end{equation}
The gravitational \emph{bulk} action has no derivatives of $K$: the only contortion dependence is quadratic. This gives us an early indication of why torsion will be algebraically determined in \emph{minimal} Einstein--Cartan theory. But that conclusion must be established by variation, and also depends on how the matter action couples to the connection.

\subsection{Why not already insert a Weyssenhoff fluid?}
A Weyssenhoff fluid is a particular model for matter carrying spin, not a property of geometry itself. Our logical order is: (1) define the general minimal gravitational action, (2) derive its equations by varying its independent fields, (3) choose a consistent spinning-matter model and determine its source terms, (4) eliminate torsion and derive effective matter corrections, and (5) specialize to a collapse metric to test whether a bounce occurs. A homogeneous bounce is not by itself proof of a baby universe or of a successful exterior matching.

The nonlinear Holst/pseudoscalaron extension belongs to a \emph{later} stage. It need not share the purely algebraic torsion dynamics of this minimal model.

\section{Technical vocabulary and distinctions}
\small
\renewcommand{\arraystretch}{1.25}
\begin{tabular}{@{}p{.25\linewidth}p{.68\linewidth}@{}}\toprule
\textbf{Term} & \textbf{Meaning in this chapter}\\\midrule
Action & Spacetime integral whose stationary variations determine equations of motion.\\
Metric $g_{\mu\nu}$ & Tensor defining spacetime intervals and index raising/lowering.\\
Connection $\Gamma^\rho{}_{\mu\nu}$ & Coefficients specifying covariant differentiation and parallel transport.\\
Metric compatibility & $\nabla_\lambda g_{\mu\nu}=0$: transported inner products are preserved.\\
Levi--Civita connection & Unique metric-compatible connection with vanishing torsion.\\
Torsion $T^\rho{}_{\mu\nu}$ & Antisymmetric part of a connection in its lower indices.\\
Contortion $K^\rho{}_{\mu\nu}$ & Tensor difference between full and Levi--Civita connections.\\
Riemann tensor & Curvature built from derivatives and products of the connection.\\
Ricci tensor / scalar & Successive contractions of the Riemann tensor.\\
Contraction & Summing over one repeated upper and lower index.\\
Covariant derivative & Tensorial differentiation using connection corrections.\\
Total divergence & A term whose invariant bulk integral is a boundary integral.\\
Independent variable & A field varied separately, subject to the chosen constraints.\\
Quadratic contortion & The specific scalar $Q(K)$ in equation \eqref{eq:QK}, not arbitrary $K^2$.\\\bottomrule
\end{tabular}
\normalsize

\section{Recap and next chapter}
We have established the following logical chain:
\begin{equation*}
\boxed{\begin{gathered}
g_{\mu\nu}\quad \longrightarrow\quad \{^\rho{}_{\mu\nu}\}(g),\\
T^\rho{}_{\mu\nu}\ \longleftrightarrow\ K^\rho{}_{\mu\nu},\qquad
\widetilde\Gamma=\{g\}+K,\\
\widetilde\Gamma\quad\longrightarrow\quad\widetilde R^\rho{}_{\sigma\mu\nu}
\quad\longrightarrow\quad\widetilde R_{\mu\nu}\quad\longrightarrow\quad\widetilde R,\\
\widetilde R=\RR+\LC_\mu(A^\mu-B^\mu)+Q(K),\\
S[g,K,\psi]=\frac1{2\kappa}\int\sqrt{-g}\,[\RR+Q(K)]\,d^4x+S_m+S_{\partial M}.
\end{gathered}}
\end{equation*}

\textbf{Chapter 2 begins by varying this action.} A metric variation gives the gravitational metric equation; a constrained contortion variation gives the Cartan spin--torsion equation. Neither variation is carried out in this chapter.

\section*{Primary sources and reading guidance}
\begin{enumerate}[leftmargin=*,itemsep=3pt]
\item S. M. Carroll, \emph{Lecture Notes on General Relativity}, \href{https://arxiv.org/abs/gr-qc/9712019}{arXiv:gr-qc/9712019}. Consult the discussion of differential geometry, curvature and the Hilbert action. Used here for the standard GR framework and motivation.
\item A. Trautman, \emph{Einstein--Cartan Theory}, \href{https://arxiv.org/abs/gr-qc/0606062}{arXiv:gr-qc/0606062}. Consult the Riemann--Cartan geometry and gravitational Lagrangian sections, especially the metricity and connection decomposition discussion. Used for the minimal EC formulation and the distinction between geometry and field equations.
\item S. D. Br\'echet, M. P. Hobson and A. N. Lasenby, \emph{Weyssenhoff fluid dynamics in general relativity using a 1+3 covariant approach}, \href{https://arxiv.org/abs/0706.2367}{arXiv:0706.2367}. See Section 2, equations (1)--(15), for their conventions and action; Section 3 for the matter model to be used later. Their torsion and contortion conventions differ from ours.
\end{enumerate}
\begin{whybox}\textbf{Scope and honesty.} The geometric identities in this chapter are mathematical consequences of our definitions and metric compatibility. The action is a physical model choice motivated by symmetry, simplicity and established weak-field behavior. A spin-induced bounce requires additional matter and dynamical assumptions and is \emph{not} derived in this chapter.\end{whybox}
\end{document}
